\documentclass[10pt,onecolumn,journal,compsoc]{IEEEtran}
\usepackage[T1]{fontenc}
\usepackage[table]{xcolor}
\usepackage{booktabs,tabularx,array,longtable}
\usepackage[nocompress]{cite}
\usepackage[colorlinks=true,allcolors=blue]{hyperref}
\definecolor{surveyheader}{HTML}{D9E2F3}
\definecolor{surveystripe}{HTML}{F3F6F9}
\renewcommand{\tabularxcolumn}[1]{>{\raggedright\arraybackslash}p{#1}}
\renewcommand{\thesection}{S\arabic{section}}
\renewcommand{\thetable}{S\arabic{table}}
\title{Supplementary Material\\Alignment in Multimodal Large Language Models: A Survey}
\author{Author information pending confirmation}
\begin{document}
\maketitle
\section{Scope and Reading Guide}
This independent supplement expands the mechanism and evaluation comparisons in the survey. It contains analytical distinctions and source locations, not new experiments or a matched model leaderboard. The Modality Gap concerns correspondence and use of task-relevant multimodal evidence; the Intent Gap concerns behavior relative to specified goals and constraints. A method can influence both, and a capability error can remain unattributed. Gap labels describe the intervention's intended target, not a claim that it closes the gap.

\section{Architecture Comparison}
Table~\ref{tab:supp-architecture} expands the evidence-access attributes compared in the main text. The relevant cost depends on input resolution, context length, trainable components, and inference budget; qualitative trade-offs should not be read as experimentally matched efficiency rankings.
\begin{table}[!ht]
\caption{Architectural interventions and the controlled comparisons needed to evaluate them.}
\label{tab:supp-architecture}
\centering\small
\setlength{\tabcolsep}{5pt}\renewcommand{\arraystretch}{1.22}
\rowcolors{2}{surveystripe}{white}
\begin{tabularx}{\textwidth}{>{\raggedright\arraybackslash}p{0.20\textwidth}XXX}
\toprule\rowcolor{surveyheader}
\textbf{Intervention / examples} & \textbf{Selection mechanism} & \textbf{Assumption to inspect} & \textbf{Controlled comparison}\\\midrule
Bridge compression: BLIP-2, TokenPacker~\cite{li2023blip2,li2025tokenpacker} & Learned query bottleneck or regional-detail injection & A compact token set retains evidence needed downstream & Match decoder and token budget; test small objects and local attributes separately from global answers.\\
Discrete interfaces: SpeechT5, AnyGPT~\cite{ao2022speecht5,zhan2024anygpt} & Codebook-based representations and token sequences & Discretization preserves task-relevant information across modalities & Separate codebook reconstruction error from reasoning and sequence-length effects.\\
Explicit structure: MAIL, VideoTree~\cite{dong2024modality,wang2025videotree} & Entity-mediated exchange or query-adaptive video hierarchy & Extraction and pruning retain decisive entities or events & Perturb or remove selected and unselected evidence; measure missed-event errors and selection cost.\\
Spatial/layer access: AGE-VLM, Dense Connector~\cite{mahajan2025attention,yao2024dense} & Spatial supervision or multi-layer feature aggregation & Additional features contribute useful detail rather than redundancy & Match image resolution and compute; isolate feature source, fusion, and spatial supervision.\\
Modality-indexed experts: VLMo, MoT~\cite{bao2022vlmo,liang2024mixture} & Modality-conditioned parameter selection with shared interaction & Specialization preserves cross-modal transfer & Report total and active parameters, per-modality loss, and held-out cross-modal tasks.\\
Adaptive experts: OneLLM, MoIIE~\cite{han2024onellm,wang2025moiie} & Learned routing among projections or experts & Routing uses task-relevant evidence rather than easy modality cues & Audit assignment distributions and evidence interventions separately; balanced utilization is not evidence reliance.\\
Graph augmentation: LLaVA-SG, MEAformer~\cite{wang2024llavasg,chen2022meaformer} & Scene relations or multimodal entity structure & Supplied relations are relevant and sufficiently correct & Compare correct, corrupted, and removed relations; distinguish the MLLM mechanism from entity-alignment foundations.\\
\bottomrule
\end{tabularx}
\end{table}

\section{Operational Diagnostic Matrix}
The contrasts in Table~\ref{tab:supp-diagnostics} test different hypotheses. Maintain task validity after an edit, hold unrelated inputs fixed, and use a reference appropriate to the edited example. A response change under an arbitrary perturbation is not by itself evidence of grounding.
\begin{table}[!ht]
\caption{Diagnostic contrasts for separating evidence failures, constraint violations, and reasoning-capability errors.}
\label{tab:supp-diagnostics}
\centering\small
\setlength{\tabcolsep}{5pt}\renewcommand{\arraystretch}{1.22}
\rowcolors{2}{surveystripe}{white}
\begin{tabularx}{\textwidth}{>{\raggedright\arraybackslash}p{0.18\textwidth}XXX}
\toprule\rowcolor{surveyheader}
\textbf{Question} & \textbf{Contrast and reference} & \textbf{Supported interpretation} & \textbf{Boundary / control}\\\midrule
Is evidence represented? & Probe or compare representations with known modality correspondence; SAIL and MIR~\cite{zhang2025assessing,huang2025deciphering}. & Evidence availability or cross-modal compatibility. & An informative representation need not be used by the decoder.\\
Is relevant evidence used? & Change a blue mug to red while keeping the color question fixed; update the correct answer. & Consistent, task-appropriate changes support image dependence. & Include irrelevant-region edits and matched natural examples to detect edit artifacts.\\
Is a claim supported? & Check object claims against annotations, as in POPE and CHAIR~\cite{li2023evaluating,rohrbach2018object}. & Unsupported-object response rate. & Does not localize the failure to encoding, fusion, or decoding.\\
Is the action constraint respected? & Keep the scene fixed and contrast ``describe without moving'' with an authorized action request. & Violations despite correct perception support an Intent Gap attribution. & Check benign utility; universal refusal is not successful task alignment.\\
Did reasoning fail after perception? & Verify the read values or objects before testing the subsequent inference; MIRAGE~\cite{dong2026mirage}. & A capability error can be distinguished from incorrect perception. & Do not force the error into either gap without a further evidence-use or constraint test.\\
Is the reward judge reliable? & Compare pairwise judgments with human-validated labels in VL-RewardBench and Multimodal RewardBench~\cite{li2025vlrewardbench,yasunaga2025multimodalrewardbench}. & Agreement on the benchmark's categories and candidates. & Audit new policy outputs, presentation order, and style confounders separately.\\
Do gains transfer? & Re-test evidence and constraint contrasts after domain or modality shifts; include uncertainty and benign-utility checks. & Transfer within the documented target conditions. & Aggregate accuracy can hide opposite movements in grounding and behavioral compliance.\\
\bottomrule
\end{tabularx}
\end{table}

\section{Related-Survey Coverage Locations}
Table~\ref{tab:supp-surveys} supports the main comparison matrix. Dedicated coverage, distributed discussion, and absence of an organizing dimension are different judgments. A cross does not mean that a paper never mentions a topic. Version-specific locations are used where a reviewed preprint differs from its later publication; these are not claims about every version.
\begin{table}[!ht]
\caption{Source locations underlying the related-survey comparison.}
\label{tab:supp-surveys}
\centering\small
\setlength{\tabcolsep}{5pt}\renewcommand{\arraystretch}{1.22}
\rowcolors{2}{surveystripe}{white}
\begin{tabularx}{\textwidth}{>{\raggedright\arraybackslash}p{0.20\textwidth}XX}
\toprule\rowcolor{surveyheader}
\textbf{Reviewed survey} & \textbf{Coverage locations} & \textbf{Interpretation for the matrix}\\\midrule
Caffagni et al.~\cite{caffagni2024revolution}, Findings ACL 2024 & Sec. 2: model components and training; Sec. 3: capabilities and other modalities; Appendices B--C: training data and evaluation benchmarks. & Objectives are discussed within training, not as a separate objective-family comparison.\\
Shu et al.~\cite{shu2025alignmentexplainability}, Findings EMNLP 2025 & Sec. 3: misalignment causes; Sec. 4: mitigation; appendices expand training stages, diagnostics, and benchmarks. & Data, model, training, and evaluation are organizing dimensions; applications are not a dedicated method grouping.\\
Lu et al.~\cite{lu2025representationpotentials}, EMNLP 2025 & Sec. 3: similarity metrics; Sec. 4: cross-domain evidence; Sec. 5: scale, architectures, objectives, and training/task diversity. & Architecture, objectives, data, training, and transfer applications are discussed, but the organizing focus is representation potential.\\
Fang et al.~\cite{fang2026alignmentsurvey}, publisher PDF & Secs. 3--7: implicit, encoder, behavioral, architectural, and inference-time alignment; Sec. 8: synthetic data and risks; Sec. 9: evaluation. & Already covers representation and behavior, inference-time intervention, data risks, and evaluation. Those topics alone are not exclusive contributions of this survey.\\
Li and Tang~\cite{li2024alignmentfusion}, arXiv v2, 11 Oct. 2025 & Sec. 2.2: datasets; Sec. 3.4: applications; Sec. 4: fusion methods; Sec. 5: challenges with training-module comparisons. & Applications are explicitly covered. Training and evaluation occur within method/challenge comparisons rather than a separate staged-training and diagnostic framework.\\
Ho et al.~\cite{ho2026hallucinationunderstanding}, publisher PDF & Sec. 4 groups mitigation by data, architecture, inference, loss optimization, and evaluation. & Substantial lifecycle coverage is acknowledged; application domains are not the central organizing structure.\\
Yu et al.~\cite{yu2025preferencesurvey}, arXiv v2, 23 Mar. 2025 & Sec. 3: preference algorithms and applications; Sec. 4: annotation; Sec. 5: evaluation, including rewards and alignment. & Preference methods, supervision, evaluation, and applications are covered; architectural families are not the organizing focus.\\
\bottomrule
\end{tabularx}
\end{table}

\section{Interpretation and Scope Boundaries}
The contribution of the main survey is the joint synthesis of intervention mechanisms, supervision assumptions, diagnostic evidence, and application constraints, not the invention of every axis in the comparison matrix. The comparisons retain foundational vision--language models when they explain an MLLM mechanism and adjacent applications when they expose a transferable constraint. Core MLLM instruction or behavioral alignment is distinguished from image--text representation learning, multimodal entity alignment, and multimodal prediction without an LLM.

Graph augmentation is treated as an architectural mechanism, with application evidence used to explain its consequences. Recommendation and climate cases illustrate supervision and evidence constraints without claiming exhaustive domain coverage. GUI-agent control and multimodal content generation are outside comprehensive coverage. Text-only preference-learning studies are methodological background, not direct demonstrations of multimodal alignment.

\clearpage
\bibliographystyle{IEEEtran}
\bibliography{references,data}
\end{document}
